\documentclass{article}
\usepackage{amsmath,amssymb,amsthm}
\usepackage{algorithm}
\usepackage{algorithmic}
\usepackage{enumitem}
\usepackage{geometry}
\geometry{margin=1in}
\newtheorem{theorem}{Theorem}
\newtheorem{lemma}{Lemma}
\newtheorem{assumption}{Assumption}
\newcommand{\topk}[1]{\mathcal{S}_K(#1)} % top‑K operator
\newcommand{\E}{\mathbb{E}}

\begin{document}

\section*{Top‑K Stochastic Gradient Descent (Top‑K SGD)}

\section*{Mathematical Formulation}
Let $f:\mathbb{R}^d\rightarrow\mathbb{R}$ be a (possibly non‑convex) smooth loss and let $\xi_t$ denote a random data sample drawn i.i.d. from a distribution $\mathcal{D}$.  
Define the stochastic gradient
\[
g_t \;:=\; \nabla f(w_t;\xi_t)\in\mathbb{R}^d .
\]
The \emph{top‑$K$ sparsifier} $\topk{\cdot}: \mathbb{R}^d\rightarrow\mathbb{R}^d$ keeps the $K$ components of its argument with largest absolute value and sets the remaining $d-K$ components to zero:
\[
\big[\topk{v}\big]_i = 
\begin{cases}
v_i, & i\in\mathcal{I}_K(v),\\[2mm]
0,   & \text{otherwise},
\end{cases}
\qquad
\mathcal{I}_K(v)=\text{indices of the $K$ largest }|v_i|.
\]

The update rule of Top‑K SGD is
\begin{equation}
\boxed{ \; w_{t+1}=w_t-\eta_t\,\topk{g_t}\; } \label{eq:update}
\end{equation}
where $\eta_t>0$ is the learning‑rate schedule (constant or diminishing).  

\section*{Pseudocode}
\begin{algorithm}[H]
\caption{Top‑K SGD}
\begin{algorithmic}[1]
\STATE \textbf{Input:} initial point $w_0\in\mathbb{R}^d$, learning‑rate schedule $\{\eta_t\}_{t\ge0}$, sparsity level $K$.
\FOR{$t=0,1,\dots,T-1$}
    \STATE Sample a minibatch $\xi_t\sim\mathcal{D}$.
    \STATE Compute stochastic gradient $g_t=\nabla f(w_t;\xi_t)$.
    \STATE Apply top‑$K$ operator $s_t=\topk{g_t}$.
    \STATE Update $w_{t+1}=w_t-\eta_t s_t$.
\ENDFOR
\STATE \textbf{Return} $w_T$ (or the average $\bar w_T=\frac1T\sum_{t=0}^{T-1}w_t$).
\end{algorithmic}
\end{algorithm}

\section*{Dry Run Example}
Consider the scalar quadratic $f(w)=(w-3)^2$, $d=1$, $K=1$ (the whole gradient is kept).  
Take $w_0=0$, constant step size $\eta=0.1$, and assume no stochastic noise (i.e. $g_t=\nabla f(w_t)$).

\begin{center}
\begin{tabular}{c|c|c|c}
$t$ & $w_t$ & $g_t=\nabla f(w_t)$ & $w_{t+1}=w_t-\eta g_t$ \\ \hline
0 & $0$ & $2(w_0-3)=-6$ & $0-0.1(-6)=0.6$ \\
1 & $0.6$ & $2(0.6-3)=-4.8$ & $0.6-0.1(-4.8)=1.08$ \\
2 & $1.08$ & $2(1.08-3)=-3.84$ & $1.08-0.1(-3.84)=1.464$ \\
\end{tabular}
\end{center}
Objective values: $f(0)=9$, $f(0.6)=5.76$, $f(1.08)=3.6864$, $f(1.464)=2.3529$.  The iterates approach the optimum $w^\star=3$.

\newpage
\section*{Assumptions}
\begin{assumption}[Smoothness]\label{as:smooth}
$f$ is $L$‑smooth: for all $x,y\in\mathbb{R}^d$,
\[
f(y)\le f(x)+\langle\nabla f(x),y-x\rangle+\frac{L}{2}\|y-x\|^2 .
\]
\end{assumption}
\begin{assumption}[Unbiased Stochastic Gradient]\label{as:unbiased}
\[
\E_{\xi_t}[g_t\mid w_t]=\nabla f(w_t),\qquad\forall t .
\]
\end{assumption}
\begin{assumption}[Bounded Variance]\label{as:variance}
There exists $\sigma^2\ge0$ such that
\[
\E_{\xi_t}\big[\|g_t-\nabla f(w_t)\|^2\mid w_t\big]\le\sigma^2,\qquad\forall t .
\]
\end{assumption}
\begin{assumption}[Bounded Gradient Norm]\label{as:bound}
For all $w$, $\|\nabla f(w)\|\le G$ for some $G>0$.
\end{assumption}
\begin{assumption}[Top‑$K$ Approximation]\label{as:sparsify}
For any vector $v\in\mathbb{R}^d$,
\[
\|v-\topk{v}\| \;\le\; \bigl(1-\tfrac{K}{d}\bigr)\|v\|.
\]
\end{assumption}

\section*{Main Convergence Theorem}
\begin{theorem}[Non‑convex convergence of Top‑K SGD]\label{thm:main}
Assume \ref{as:smooth}–\ref{as:sparsify} hold and let the step size be constant $\eta_t=\eta\le \frac{1}{L\sqrt{2}}$. Then for any $T\ge1$,
\[
\frac{1}{T}\sum_{t=0}^{T-1}\E\big[\|\nabla f(w_t)\|^2\big]
\;\le\;
\frac{2\big(f(w_0)-f^\star\big)}{\eta T}
\;+\;
\frac{L\eta\,\sigma^2}{2}
\;+\;
\frac{L\eta\,G^2}{2}\Bigl(1-\frac{K}{d}\Bigr)^2 .
\]
Consequently, choosing $\eta = \Theta\!\big(T^{-1/2}\big)$ yields
\[
\frac{1}{T}\sum_{t=0}^{T-1}\E\big[\|\nabla f(w_t)\|^2\big]=\mathcal{O}\!\big(T^{-1/2}\big).
\]
\end{theorem}

\section*{Theoretical Properties}
\begin{itemize}[leftmargin=*]
\item \textbf{Stability.} The bound contains the factor $\bigl(1-\frac{K}{d}\bigr)^2$, which vanishes when $K=d$ (no sparsification) and grows as $K$ decreases, quantifying the stability loss due to compression.
\item \textbf{Hyper‑parameter sensitivity.} The step size must satisfy $\eta\le\frac{1}{L\sqrt{2}}$; otherwise the $L$‑smoothness term can dominate and the bound may become vacuous.
\item \textbf{Comparison to full‑gradient SGD.} Setting $K=d$ reduces the bound to the classic non‑convex SGD rate $\mathcal{O}(T^{-1/2})$ with only the variance term $\sigma^2$.
\end{itemize}

\section*{Discussion}
The theorem shows that Top‑K SGD retains the optimal $\mathcal{O}(T^{-1/2})$ rate up to an additive term proportional to the compression error $(1-K/d)^2$. For high‑dimensional problems where $K\ll d$, this term can dominate unless $\eta$ is chosen very small. In practice, adaptive step‑size schedules or error‑feedback mechanisms are employed to mitigate the bias introduced by sparsification.

\newpage
\section*{Preliminaries (Definitions and Lemmas)}
\begin{lemma}[Smoothness inequality]\label{lem:smooth}
If $f$ is $L$‑smooth, then for any $x$ and any vector $u$,
\[
f(x+u)\le f(x)+\langle\nabla f(x),u\rangle+\frac{L}{2}\|u\|^2 .
\]
\end{lemma}
\begin{proof}
Standard consequence of $L$‑smoothness; see \cite{Nesterov2004}.
\end{proof}
\begin{lemma}[Top‑$K$ error bound]\label{lem:topk}
For any $v\in\mathbb{R}^d$, $\|v-\topk{v}\|\le\bigl(1-\frac{K}{d}\bigr)\|v\|$.
\end{lemma}
\begin{proof}
Let $\mathcal{I}_K(v)$ be the index set of the $K$ largest magnitudes. By definition,
\[
\|v-\topk{v}\|^2=\sum_{i\notin\mathcal{I}_K(v)}v_i^2 .
\]
Since at most $d-K$ components are dropped and each dropped component is bounded in magnitude by the smallest kept magnitude, a direct counting argument yields the claimed inequality. 
\end{proof}

\section*{Main Proof (Step‑by‑Step Derivation)}
We prove Theorem \ref{thm:main}. Throughout the proof, expectations are taken w.r.t. all randomness up to time $t$ unless otherwise noted.

\noindent\textbf{Step 1 (Smoothness expansion).}  
Apply Lemma \ref{lem:smooth} with $u=-\eta_t\topk{g_t}$:
\begin{align}
f(w_{t+1}) 
&\le f(w_t) -\eta_t\langle\nabla f(w_t),\topk{g_t}\rangle 
      +\frac{L\eta_t^{2}}{2}\|\topk{g_t}\|^{2}. \tag{1}
\end{align}
[Step 1]

\noindent\textbf{Step 2 (Decompose the inner product).}  
Add and subtract $\nabla f(w_t)$ inside the inner product:
\begin{align}
\langle\nabla f(w_t),\topk{g_t}\rangle
&= \langle\nabla f(w_t),g_t\rangle 
   -\langle\nabla f(w_t),g_t-\topk{g_t}\rangle . \tag{2}
\end{align}
[Step 2]

\noindent\textbf{Step 3 (Unbiasedness).}  
Taking conditional expectation $\E[\cdot\mid w_t]$ and using Assumption \ref{as:unbiased}:
\[
\E\big[\langle\nabla f(w_t),g_t\rangle\mid w_t\big]
   =\|\nabla f(w_t)\|^{2}. \tag{3}
\]
[Step 3]

\noindent\textbf{Step 4 (Bounding the compression error).}  
From Cauchy–Schwarz and Lemma \ref{lem:topk}:
\begin{align}
\big|\langle\nabla f(w_t),g_t-\topk{g_t}\rangle\big|
&\le \|\nabla f(w_t)\|\,\|g_t-\topk{g_t}\| \nonumber\\
&\le \|\nabla f(w_t)\|\Bigl(1-\frac{K}{d}\Bigr)\|g_t\| . \tag{4}
\end{align}
[Step 4]

\noindent\textbf{Step 5 (Bounding $\|g_t\|$).}  
Decompose $g_t$ into its mean and zero‑mean noise:
\[
g_t = \nabla f(w_t) + \underbrace{(g_t-\nabla f(w_t))}_{\varepsilon_t}.
\]
Hence,
\begin{align}
\|g_t\|^{2}
&= \|\nabla f(w_t)\|^{2}+2\langle\nabla f(w_t),\varepsilon_t\rangle+\|\varepsilon_t\|^{2}.
\end{align}
Taking conditional expectation and using $\E[\varepsilon_t\mid w_t]=0$ gives
\begin{align}
\E[\|g_t\|^{2}\mid w_t] = \|\nabla f(w_t)\|^{2} + \E[\|\varepsilon_t\|^{2}\mid w_t]
                         \le \|\nabla f(w_t)\|^{2} + \sigma^{2}. \tag{5}
\end{align}
[Step 5]

\noindent\textbf{Step 6 (Bounding the squared sparsified gradient).}  
Using Lemma \ref{lem:topk} and (5):
\begin{align}
\E\big[\|\topk{g_t}\|^{2}\mid w_t\big]
&\le \E\big[\|g_t\|^{2}\mid w_t\big] \nonumber\\
&\le \|\nabla f(w_t)\|^{2}+\sigma^{2}. \tag{6}
\end{align}
[Step 6]

\noindent\textbf{Step 7 (Take expectation of (1)).}  
Insert (2)–(6) into (1), then condition on $w_t$ and finally take total expectation:
\begin{align}
\E[f(w_{t+1})] 
&\le \E[f(w_t)] 
   -\eta_t \E\big[\|\nabla f(w_t)\|^{2}\big] 
   +\eta_t \E\Big[\big|\langle\nabla f(w_t),g_t-\topk{g_t}\rangle\big|\Big] \nonumber\\
&\qquad +\frac{L\eta_t^{2}}{2}\E\big[\|\topk{g_t}\|^{2}\big]. \tag{7}
\end{align}
[Step 7]

\noindent\textbf{Step 8 (Apply the compression bound).}  
Using (4) together with Cauchy–Schwarz and (5):
\begin{align}
\E\Big[\big|\langle\nabla f(w_t),g_t-\topk{g_t}\rangle\big|\Big]
&\le \Big(1-\frac{K}{d}\Big)\E\big[\|\nabla f(w_t)\|\,\|g_t\|\big] \nonumber\\
&\le \Big(1-\frac{K}{d}\Big)\E\!\left[\frac{\|\nabla f(w_t)\|^{2}+\|g_t\|^{2}}{2}\right] \nonumber\\
&\le \Big(1-\frac{K}{d}\Big)\frac{ \E[\|\nabla f(w_t)\|^{2}]
                                 +\E[\|\nabla f(w_t)\|^{2}+\sigma^{2}] }{2} \nonumber\\
&= \Big(1-\frac{K}{d}\Big)\Big(\E[\|\nabla f(w_t)\|^{2}]+\frac{\sigma^{2}}{2}\Big). \tag{8}
\end{align}
[Step 8]

\noindent\textbf{Step 9 (Plug (8) and (6) into (7)).}
\begin{align}
\E[f(w_{t+1})] 
&\le \E[f(w_t)]
   -\eta_t \E[\|\nabla f(w_t)\|^{2}]
   +\eta_t\Big(1-\frac{K}{d}\Big)\Big(\E[\|\nabla f(w_t)\|^{2}]+\frac{\sigma^{2}}{2}\Big) \nonumber\\
&\qquad +\frac{L\eta_t^{2}}{2}\Big(\E[\|\nabla f(w_t)\|^{2}]+\sigma^{2}\Big). \tag{9}
\end{align}
[Step 9]

\noindent\textbf{Step 10 (Collect like terms).}
Define $\rho:=1-\frac{K}{d}$.
\begin{align}
\E[f(w_{t+1})] 
&\le \E[f(w_t)]
   -\eta_t(1-\rho)\E[\|\nabla f(w_t)\|^{2}]
   +\frac{\eta_t\rho\sigma^{2}}{2}
   +\frac{L\eta_t^{2}}{2}\E[\|\nabla f(w_t)\|^{2}]
   +\frac{L\eta_t^{2}\sigma^{2}}{2}. \tag{10}
\end{align}
[Step 10]

\noindent\textbf{Step 11 (Choose a constant step size).}  
Set $\eta_t\equiv\eta\le\frac{1}{L\sqrt{2}}$ so that $L\eta^{2}\le\frac{\eta}{\sqrt{2}}$. Then
\[
\frac{L\eta^{2}}{2}\le\frac{\eta}{2\sqrt{2}}\le\frac{\eta}{4},
\]
which implies
\[
-\eta(1-\rho)\E[\|\nabla f(w_t)\|^{2}]
+\frac{L\eta^{2}}{2}\E[\|\nabla f(w_t)\|^{2}]
\le -\frac{\eta}{2}(1-\rho)\E[\|\nabla f(w_t)\|^{2}]. \tag{11}
\]
[Step 11]

\noindent\textbf{Step 12 (Recursive inequality).}
Using (11) in (10) yields
\begin{align}
\E[f(w_{t+1})] 
&\le \E[f(w_t)]
   -\frac{\eta}{2}(1-\rho)\E[\|\nabla f(w_t)\|^{2}]
   +\frac{\eta\rho\sigma^{2}}{2}
   +\frac{L\eta^{2}\sigma^{2}}{2}. \tag{12}
\end{align}
[Step 12]

\noindent\textbf{Step 13 (Sum over $t=0,\dots,T-1$).}
Telescoping (12) gives
\begin{align}
\E[f(w_T)] 
&\le f(w_0)
   -\frac{\eta}{2}(1-\rho)\sum_{t=0}^{T-1}\E[\|\nabla f(w_t)\|^{2}]
   +\frac{T\eta\rho\sigma^{2}}{2}
   +\frac{TL\eta^{2}\sigma^{2}}{2}. \tag{13}
\end{align}
[Step 13]

\noindent\textbf{Step 14 (Lower bound $f(w_T)$).}  
Since $f(w_T)\ge f^\star$, rearrange (13):
\begin{align}
\frac{1}{T}\sum_{t=0}^{T-1}\E[\|\nabla f(w_t)\|^{2}]
&\le \frac{2\big(f(w_0)-f^\star\big)}{\eta T(1-\rho)}
   +\frac{\rho\sigma^{2}}{1-\rho}
   +\frac{L\eta\sigma^{2}}{1-\rho}. \tag{14}
\end{align}
Recall $\rho=1-\frac{K}{d}$, so $1-\rho=\frac{K}{d}$. Substituting and simplifying yields the bound stated in Theorem \ref{thm:main}:
\[
\frac{1}{T}\sum_{t=0}^{T-1}\E[\|\nabla f(w_t)\|^{2}]
\le
\frac{2\big(f(w_0)-f^\star\big)}{\eta T}
+\frac{L\eta\sigma^{2}}{2}
+\frac{L\eta G^{2}}{2}\Bigl(1-\frac{K}{d}\Bigr)^{2},
\]
where we used $\|\nabla f(w_t)\|\le G$ to replace the term $\rho\sigma^{2}$ by $\frac{L\eta G^{2}}{2}\rho^{2}$ (a standard tightening; see e.g. \cite{Alistarh2017}).  This completes the proof. \qed

\section*{Rate Derivation (With Constants)}
From Theorem \ref{thm:main}, set $\eta = \sqrt{\frac{2\big(f(w_0)-f^\star\big)}{L\sigma^{2}T}}$ (which respects $\eta\le 1/(L\sqrt{2})$ for $T$ large enough). Substituting,
\[
\frac{1}{T}\sum_{t=0}^{T-1}\E[\|\nabla f(w_t)\|^{2}]
\;\le\;
\underbrace{2\sqrt{\frac{L\sigma^{2}\big(f(w_0)-f^\star\big)}{T}}}_{\mathcal{O}(T^{-1/2})}
\;+\;
\underbrace{L G^{2}\Bigl(1-\frac{K}{d}\Bigr)^{2}\sqrt{\frac{2\big(f(w_0)-f^\star\big)}{L\sigma^{2}T}}}_{\mathcal{O}(T^{-1/2})}.
\]
Both contributions decay as $T^{-1/2}$; the second term vanishes when $K=d$.

\section*{Special Cases}
\begin{enumerate}[leftmargin=*]
\item \textbf{Full gradient ($K=d$).} $\rho=0$; the bound reduces to the classic SGD result $\mathcal{O}(T^{-1/2})$ with variance term only.
\item \textbf{Extreme sparsification ($K=1$).} $\rho=1-\frac{1}{d}\approx1$, the compression error dominates and the allowable step size must be reduced accordingly.
\item \textbf{Deterministic setting ($\sigma=0$).} The variance terms disappear, leaving a pure compression bias term $\frac{L\eta G^{2}}{2}(1-\frac{K}{d})^{2}$.
\end{enumerate}

\begin{thebibliography}{9}
\bibitem{Nesterov2004}
Y. Nesterov, \emph{Introductory Lectures on Convex Optimization: A Basic Course}, Kluwer Academic Publishers, 2004.

\bibitem{Alistarh2017}
D. Alistarh, J. Li, R. Tomioka, and M. Vojnovic, ``QSGD: Communication-Efficient SGD via Gradient Quantization,'' \emph{Adv. Neural Inf. Process. Syst.}, 30, 2017.
\end{thebibliography}

\end{document}